% ============================================================================
\mainsection{บทนำ (Introduction)}
% ============================================================================

โครงงาน \textenglish{Slime Maze} คือโครงงานการประยุกต์ใช้ความรู้ในรายวิชา \textenglish{CPE222} อิเล็กทรอนิกส์ดิจิทัลและการออกแบบวงจรเชิงตรรกะ (Digital Electronics and Logic Design) มาออกแบบและพัฒนาชิ้นงานจริงบนฮาร์ดแวร์ FPGA โดยคณะผู้จัดทำมีความสนใจในการพัฒนาระบบเกมอาร์เคดที่มีแนวคิดต้นแบบมาจากเกมแนวเขาวงกตไล่ล่า (Maze-Chase Arcade Game) อย่าง \textenglish{Pac-Man} ผสานกับระบบการวิวัฒนาการของตัวละครสไลม์ 3 ขั้น (Slime Evolution) ที่พัฒนาด้วยภาษา Verilog HDL และแสดงผลผ่านพอร์ต VGA ความละเอียด $640\times480$ พิกเซล บนบอร์ด Digilent Nexys A7-100T โดยมีรายละเอียดพื้นฐานที่เกี่ยวข้องดังนี้

\begin{figure}[H]
    \centering
    \smartimage{0.82\textwidth}{banner.png}{4.0cm}{ภาพแบนเนอร์แนะนำเกม Slime Maze (images/banner.png)}
    \caption{ภาพแสดงผลหน้าปกและแนวคิดของเกม Slime Maze}
\end{figure}

\subsection{แนวคิดและแรงบันดาลใจของเกมเขาวงกต (Arcade Maze Game)}

เกมแนวเขาวงกตไล่ล่า (Maze-Chase Game) เป็นเกมอาร์เคดคลาสสิกที่มีต้นแบบสำคัญมาจากเกม Pac-Man โดยผู้เล่นจะต้องบังคับตัวละครเคลื่อนที่ไปตามทางเดินในเขาวงกตแบบ 2 มิติ เพื่อเก็บเม็ดคะแนน (Pellets) ให้ครบตามเงื่อนไข พร้อมกับหลบหลีกศัตรู (Ghosts) ที่คอยไล่ล่าจากทิศทางต่างๆ แต่เมื่อเก็บไอเทมพิเศษได้ บทบาทจะพลิกกลับจาก “ผู้ถูกล่า” เป็น “ผู้ล่า” ทำให้สามารถไล่กินศัตรูเพื่อทำคะแนนโบนัสได้ชั่วคราว

โครงงาน Slime Maze ได้นำแนวคิดดังกล่าวมายกระดับด้วยระบบวิวัฒนาการตัวละคร 3 ร่าง (Baby Slime $\rightarrow$ Big Slime $\rightarrow$ King Slime) โดยผู้เล่นต้องเก็บเม็ดพลังงานในเขาวงกตขนาด $20\times15$ ช่อง ให้ครบ 30 เม็ดเพื่อเติบโตจากร่างที่ 1 เป็นร่างที่ 2 จากนั้นมงกุฎ (Crown) จะปรากฏขึ้นตามมุมของแผนที่ เมื่อเก็บมงกุฎได้ สไลม์จะเข้าสู่ “โหมดราชา” (King Mode / ร่างที่ 3) เป็นเวลา 10 วินาที ซึ่งจะเดินเร็วขึ้น 2 เท่า และสามารถไล่กินผีทั้ง 4 ตัวเพื่อเอาชนะเกมได้ หากพลาดถูกผีจับตัวครบ 3 ครั้ง เกมจะจบลงทันที

% [INSERT PICTURE: pacman_ref.png — ภาพหน้าจอเกม Pac-Man ต้นแบบเพื่อเปรียบเทียบแรงบันดาลใจ]
\begin{figure}[H]
    \centering
    \smartimage{0.68\textwidth}{pacman_ref.png}{4.2cm}{ภาพหน้าจอ (Screenshot) ของเกม Pac-Man ต้นแบบ แสดงแผนที่เขาวงกต เม็ดคะแนน ตัวละครผู้เล่น และผี 4 ตัว เพื่อใช้อ้างอิงแรงบันดาลใจในการออกแบบ (เพิ่ม URL ที่มาในบรรณานุกรม)}
    \caption{ตัวอย่างภาพแสดงผลเกมต้นแบบแนวเขาวงกตไล่ล่า (Pac-Man)}
\end{figure}

\subsection{Field Programmable Gate Array (FPGA)}

FPGA เป็นอุปกรณ์วงจรรวม (Integrated Circuit) ที่ผู้ใช้สามารถเขียนโปรแกรมเพื่อกำหนดโครงสร้างวงจรดิจิทัลภายในได้ตามต้องการ ภายในชิป FPGA ประกอบด้วยหน่วยตรรกะพื้นฐานจำนวนมาก เช่น Look-Up Tables (LUTs), Flip-Flops, Block RAM (BRAM) และตัวจัดการสัญญาณนาฬิกา (Clock Management Tiles / MMCM) ที่สามารถเชื่อมต่อกันเป็นวงจรดิจิทัลเชิงตรรกะและวงจรลำดับได้อย่างหลากหลาย เช่น วงจรนับ (Counter), วงจรบวก, เครื่องสถานะจำกัด (Finite State Machine: FSM), หน่วยความจำแบบ ROM และวงจรสร้างสัญญาณภาพ (Video Controller)

จุดเด่นสำคัญของ FPGA คือความสามารถในการประมวลผลแบบคู่ขนาน (Parallel Processing) ในระดับฮาร์ดแวร์โดยไม่ต้องพึ่งพาระบบปฏิบัติการหรือซีพียูแบบซอฟต์แวร์ การพัฒนาวงจรบน FPGA ใช้ภาษาคำอธิบายฮาร์ดแวร์ (Hardware Description Language: HDL) เช่น Verilog หรือ VHDL จากนั้นจึงนำโค้ดผ่านกระบวนการ Synthesis และ Implementation เพื่อสร้างไฟล์กำหนดค่าที่เรียกว่า Bitstream สำหรับโปรแกรมลงบนชิป FPGA ต่อไป

\subsection{Digilent Nexys A7-100T Artix-7 FPGA Trainer Board}

บอร์ดพัฒนา Digilent Nexys A7-100T เป็นบอร์ดฝึกปฏิบัติและพัฒนาวงจรดิจิทัลที่มาพร้อมกับชิป FPGA ตระกูล Xilinx Artix-7 รหัส \texttt{xc7a100tcsg324-1} ซึ่งเป็นรุ่นที่มีความจุทางตรรกะสูงเพียงพอต่อการออกแบบระบบเกมและกราฟิกแบบเต็มรูปแบบ บนตัวบอร์ดติดตั้งอุปกรณ์ต่อพ่วงพื้นฐานมาครบถ้วนพร้อมใช้งานทันทีโดยไม่ต้องต่อวงจรภายนอกเพิ่มเติม ได้แก่

\begin{itemize}[leftmargin=1.5cm]
    \item สัญญาณนาฬิกาหลักบนบอร์ด (On-board Oscillator) ความถี่ 100 MHz (ขา E3)
    \item ปุ่มกดแบบ Pushbutton 5 ทิศทาง (\texttt{BTNU}, \texttt{BTND}, \texttt{BTNL}, \texttt{BTNR}, \texttt{BTNC}) และปุ่มรีเซ็ต \texttt{CPU\_RESETN} (Active-Low)
    \item หลอดไฟ LED แสดงสถานะจำนวน 16 ดวง (\texttt{LED[15:0]})
    \item จอแสดงผลตัวเลข 7-Segment แบบ Common Anode จำนวน 8 หลัก (\texttt{SEG[6:0]}, \texttt{DP}, \texttt{AN[7:0]})
    \item พอร์ตเชื่อมต่อสัญญาณภาพ VGA แบบ 12-bit Color (4-bit Red, 4-bit Green, 4-bit Blue พร้อมสัญญาณ \texttt{VGA\_HS} และ \texttt{VGA\_VS}) ผ่านวงจร R-2R Ladder DAC บนบอร์ด
    \item พอร์ต Micro-USB สำหรับจ่ายไฟเลี้ยงและโปรแกรม Bitstream ผ่านวงจร USB-JTAG
\end{itemize}
